% AUTOGENERATED DRAFT from manuscript source: manuscript/appendices/A-symbols.md
% Converter: scripts/convert_markdown_chapter_to_tex.py
% Review gate: docs/latex-migration-plan.md (Phase 2 per-chapter gate)

\chapter{符号、时间层与源码变量}
\label{chap:A-symbols}

本附录服务于正文阅读和源码检索。公式中的符号采用连续模型或离散模型的常用记号；反引号中的名称则优先对应 WarpX 源码、输入文件或 diagnostics 输出中的实际字段。除特别说明外，SI 制单位适用，粒子权重 \cpath{w\_p} 表示一个宏粒子代表的真实粒子数。

使用本附录时，先区分三件事：数学符号描述的物理量、WarpX 内部数组实际保存的量、输入文件或 diagnostics 对同一名称采用的单位。它们常常有关联，但不必完全相同；尤其是 \code{ux, uy, uz} 不能只按变量名猜测单位。

\section{A.1 连续模型符号}

% <-- table block (source: A-symbols L9)
\begin{tabularx}{\textwidth}{LLL}
\toprule
符号 & 含义 & 常见单位或类型 \\
\midrule
\endfirsthead
\toprule
符号 & 含义 & 常见单位或类型 \\
\midrule
\endhead
$s$ & species 索引，例如 electron、ion、photon & 无量纲整数 \\
$f_s(\mathbf{x},\mathbf{p},t)$ & species $s$ 的相空间分布函数 & 依定义而定 \\
$q_s,m_s$ & species 电荷和静质量 & C、kg \\
$\mathbf{x}=(x,y,z)$ & 物理空间位置 & m \\
$\mathbf{p}$ & 物理动量 & kg m s$^{-1}$ \\
$\boldsymbol{\beta}=\mathbf{v}/c$ & 无量纲速度 & 无量纲 \\
$\mathbf{u}=\mathbf{p}/m=\gamma\mathbf{v}$ & proper velocity；WarpX 粒子容器内部的 \code{ux, uy, uz} 以此量保存 & m s$^{-1}$ \\
$\mathbf{p}/(mc)=\mathbf{u}/c=\gamma\boldsymbol{\beta}$ & 无量纲归一化动量；多数 species 输入的 \cpath{u} 参数采用这一约定 & 无量纲 \\
$\mathbf{v}$ & 粒子速度 & m s$^{-1}$ \\
$\gamma$ & Lorentz 因子，$\gamma=(1-v^2/c^2)^{-1/2}$ & 无量纲 \\
$c$ & 真空光速 & m s$^{-1}$ \\
$\mathbf{E},\mathbf{B}$ & 电场和磁场 & V m$^{-1}$、T \\
$\rho,\mathbf{J}$ & 电荷密度和电流密度 & C m$^{-3}$、A m$^{-2}$ \\
$\epsilon_0,\mu_0$ & 真空介电常数和磁导率 & SI 常数 \\
$w_p$ & 宏粒子权重，代表的真实粒子数 & 无量纲或按模型定义 \\
$S$ & 粒子到网格的空间形函数 & 按离散归一化定义 \\
\bottomrule
\end{tabularx}

\subsection{\texttt{ux, uy, uz} 的两层约定}

对有质量粒子，三种写法的关系是

\[
\mathbf{p}=m\mathbf{u},\qquad
\mathbf{u}=\gamma\mathbf{v},\qquad
\frac{\mathbf{u}}{c}=\frac{\mathbf{p}}{mc}=\gamma\boldsymbol{\beta},
\]

并且

\[
\gamma=\sqrt{1+\left|\frac{\mathbf{u}}{c}\right|^2}.
\]

这解释了两个表面矛盾、实际一致的源码约定。

\begin{enumerate}
\item \cpath{Source/Particles/WarpXParticleContainer.H} 将粒子容器中的 \code{ux, uy, uz} 说明为 proper velocity，即 \cpath{gamma*v}，内部量纲为 m/s。
\item \cpath{Docs/source/usage/parameters.rst} 将 \cpath{single\_particle\_u}、\cpath{multiple\_particles\_u*} 和常见 species momentum distribution 的输入定义为 \cpath{gamma*beta}，即无量纲的 \cpath{gamma*v/c}。初始化时会乘以 \cpath{c}，再写入内部粒子数组；\cpath{PlasmaInjector.cpp} 中的 \cpath{multiple\_particles\_ux} 转换正是这一边界的例子。
\end{enumerate}

因此，阅读输入文件时把 \code{uz = 10} 理解为 $\gamma\beta_z=10$；阅读粒子容器或以 SI 单位计算能量的 kernel 时，把内部 \cpath{uz} 理解为 $\gamma v_z=10c$。诊断 parser、openPMD 记录和不同输出格式还可能按其文档以 $\gamma v/c$ 暴露动量分量；遇到 \cpath{ux} 时，先确认它来自输入、容器还是输出 metadata，再进行单位换算。

\section{A.2 网格、时间层与离散量}

% <-- table block (source: A-symbols L53)
\begin{tabularx}{\textwidth}{LLL}
\toprule
符号 & 含义 & 在程序中的对应 \\
\midrule
\endfirsthead
\toprule
符号 & 含义 & 在程序中的对应 \\
\midrule
\endhead
$t^n$ & 第 $n$ 个整数时间层 & \cpath{step}、\cpath{istep} \\
$t^{n+1/2}$ & 半时间层，常用于 leapfrog 电流或磁场 & \codeesc{J\textasciicircum{}\{n+1/2\}} 等时间层语义 \\
$\Delta t$ & 时间步长 & \cpath{dt}、\cpath{warpx.const\_dt} 或派生时间步长 \\
$\Delta x,\Delta y,\Delta z$ & 各方向网格间距 & \cpath{amr.n\_cell}、几何 \cpath{cell\_size} \\
$i,j,k$ & 网格单元或节点索引 & AMReX \cpath{IntVect} 分量 \\
$\ell$ & AMR level 索引，$\ell=0$ 为最粗层 & \cpath{lev}、\cpath{level} \\
$r$ & RZ/cylindrical 几何中的径向坐标 & m \\
\cpath{rank} 或 $r_\mathrm{MPI}$ & MPI rank 或 rank 索引；不用裸 $r$ 与径向坐标混写 & \cpath{ParallelDescriptor::MyProc()} 等 \\
$V_i$ & 网格单元体积 & Cartesian 中常为 $\Delta x\Delta y\Delta z$ \\
$\rho_i^n$ & 单元/节点 $i$ 在时间层 $n$ 的电荷密度 & \cpath{rho}、\cpath{rho\_fp}、\cpath{rho\_buf} \\
$\mathbf{J}_i^{n+1/2}$ & 时间层 $n+1/2$ 的电流密度 & \cpath{current\_fp}、\cpath{current\_buf} \\
$\nabla_h\cdot\mathbf{J}$ & 离散散度 & 由 staggering 和差分 stencil 决定 \\
$S_i(\mathbf{x}_p)$ & 粒子位置对网格量 $i$ 的形函数值 & \cpath{ShapeFactors} 中的 shape 权重 \\
\bottomrule
\end{tabularx}

在 AMR 语境中，\cpath{\_fp} 与 \cpath{\_cp} 通常分别标记 fine patch 和 coarse patch，\cpath{\_aux} 则标记供粒子 gather 使用的辅助场。它们不是同一物理场的任意别名：粗细网格覆盖、guard cell 和同步阶段决定某个数组何时有效。第 3、5 和 7 章中的 \cpath{current\_fp/current\_buf}、\cpath{rho\_fp/rho\_buf} 与 \cpath{Efield\_aux/Bfield\_aux} 都应按这个生命周期来读。

\section{A.3 沉积、推进与场求解记号}

% <-- table block (source: A-symbols L73)
\begin{tabularx}{\textwidth}{LLL}
\toprule
符号或名称 & 含义 & 阅读提示 \\
\midrule
\endfirsthead
\toprule
符号或名称 & 含义 & 阅读提示 \\
\midrule
\endhead
$\rho^n\rightarrow\mathbf{J}^{n+1/2}\rightarrow\rho^{n+1}$ & 电荷守恒 current-deposition 主线的时间层关系 & 不能把 charge deposition 和 current deposition 当成同一个 kernel \\
\cpath{old} / \cpath{new} & 粒子在一个推进步或 segment 两端的形函数/位置状态 & Esirkepov、Villasenor kernel 中常见 \\
\cpath{relative\_time} & 在同一步内取样粒子位置的相对时间参数 & 影响 source 的时间层，而不是额外推进粒子 \\
\cpath{icomp} & charge component 或时间层分量选择 & 需结合调用入口解释，不能只按名字猜测 \\
\cpath{depos\_lev} & charge 沉积目标 AMR level & coarse-buffer 粒子可能直接沉积到 \cpath{lev-1} \\
\cpath{current\_fp} & fine-patch current buffer & 主要对应 fine patch 粒子沉积 \\
\cpath{current\_buf} & coarse-buffer current buffer & 后续还要经过同步/整理 \\
\cpath{rho\_fp} / \cpath{rho\_buf} & fine-patch / coarse-buffer charge buffer & 是 source route 的观测面，不等于最终 plotfile 字段 \\
\cpath{Efield\_fp} / \cpath{Bfield\_fp} & fine-patch 电场/磁场 & 不能直接假定就是粒子 gather 的数组 \\
\cpath{Efield\_cp} / \cpath{Bfield\_cp} & coarse-patch 电场/磁场 & 在 mesh refinement 近边界的 gather/sync 语义要结合 level 与 patch type 判断 \\
\cpath{Efield\_aux} / \cpath{Bfield\_aux} & 供粒子 gather 的辅助场 & 由 \cpath{UpdateAuxiliaryData} 等路径从 patch 场构造；只定义 E/B，不把它误认为 current 或 charge 容器 \\
\cpath{Direct} & 直接速度加权电流沉积 & 简单但不自动满足离散连续性方程 \\
\cpath{Esirkepov} & 基于轨迹与形函数差分的 charge-conserving current deposition & 重点看 density decomposition 和 prefix accumulation \\
\cpath{Villasenor} & 基于 cell crossing segment 的 charge-conserving current deposition & 重点看 crossing、segment 和局部 face flux \\
\cpath{Vay} & Vay current deposition / spectral 相关路径 & 常与 PSATD、current correction 组合讨论 \\
\cpath{FDTD} & 有限差分时域场求解器 & 依赖 staggered grid 和 CFL 限制 \\
\cpath{PSATD} & pseudo-spectral analytical time-domain 求解器 & 重点看谱空间系数、源项时间模型和周期边界假设 \\
\cpath{JRhom} & PSATD 的多次 $J/\rho$ 时间采样路径 & 不是多次粒子 push，而是同一轨迹的多时刻源项 \\
\cpath{PML} & 吸收边界层 & 通过 split-field 或等价谱推进吸收出射波 \\
\bottomrule
\end{tabularx}

\section{A.4 诊断、文件与验证术语}

% <-- table block (source: A-symbols L97)
\begin{tabularx}{\textwidth}{LL}
\toprule
名称 & 含义 \\
\midrule
\endfirsthead
\toprule
名称 & 含义 \\
\midrule
\endhead
plotfile & 分析用的网格/粒子输出；不可 restart，重启使用 \cpath{Full} checkpoint \\
openPMD & 粒子、网格或 reduced diagnostics 的结构化输出接口 \\
\cpath{diagNNNNNNN} & diagnostics 输出的迭代目录，例如 \cpath{diag1000000} \\
\cpath{reduced\_diags} & 不保存完整场/粒子，而输出聚合标量或低维量的 diagnostics \\
\cpath{FieldProbe} & 沿指定几何路径采样场并做积分或线探针输出的诊断 \\
\cpath{BoundaryScrapingDiagnostics} & 记录穿过粒子边界的粒子及其统计量 \\
producer & 产生字段、粒子或 reduced output 的运行时路径 \\
consumer & 读取产物并执行物理、格式或性能断言的分析脚本 \\
physics gate & 对解析解、守恒量、谱或物理量的数值容差断言 \\
writer gate & 对文件、字段、维度、iteration、轴和有限非零数据的断言 \\
checksum gate & 对最终输出做确定性校验；不能自动等价为物理正确性 \\
reader-side analysis & 从已有 plotfile/openPMD 重新读取数据并验证合同的分析层 \\
\bottomrule
\end{tabularx}

\section{A.5 常用缩写}

% <-- table block (source: A-symbols L114)
\begin{tabularx}{\textwidth}{LLL}
\toprule
缩写 & 全称 & 本书中的语境 \\
\midrule
\endfirsthead
\toprule
缩写 & 全称 & 本书中的语境 \\
\midrule
\endhead
PIC & Particle-In-Cell & 粒子-网格数值方法 \\
AMR & Adaptive Mesh Refinement & 多层网格和 coarse/fine 同步 \\
CFL & Courant-Friedrichs-Lewy & 显式场推进稳定性限制 \\
NCI & Numerical Cherenkov Instability & boosted-frame / PSATD 中的数值不稳定性 \\
RZ & cylindrical geometry with azimuthal modes & WarpX 的轴对称/模态几何 \\
EB & Embedded Boundary & 嵌入边界和 cut-cell 几何 \\
MPI & Message Passing Interface & rank 间并行通信 \\
GPU & Graphics Processing Unit & device kernel 和 GPU memory 路径 \\
QED & Quantum Electrodynamics & 高场量子辐射和 pair-production 模型 \\
LWFA / PWFA & Laser / Plasma Wakefield Acceleration & 激光/束流驱动尾场应用 \\
\bottomrule
\end{tabularx}

\section{A.6 使用规则}

\begin{enumerate}
\item 看到 \code{ux, uy, uz} 时，先确认它来自输入、内部粒子容器还是 diagnostics；分别按 $\gamma\beta$、$\gamma\mathbf{v}$ 或该输出的 metadata 解释，不能把数值直接混用。
\item 看到 \cpath{rho}、\cpath{current\_fp} 或 \cpath{current\_buf} 时，先判断它属于 local kernel、level buffer、同步后场，还是最终 diagnostics 输出；同名物理量可能处于不同生命周期阶段。
\item 看到上标 $n$、$n+1/2$ 或 \cpath{relative\_time} 时，先确认粒子位置、场、current 和 charge 是否处在同一个时间层；不能只根据变量名推断守恒性。
\end{enumerate}
